% !TeX program = xelatex
% Compile with: latexmk main.tex   (XeLaTeX and Biber, set up in latexmkrc)

% Class options: aspectratio=169 for 16:9 slides, t to top-align every frame
% (frames are vertically centred by default), handout for a printable copy.
\documentclass[10pt]{beamer}

% ---------------------------------------------------------------------------- %
%   Preamble                                                                   %
% ---------------------------------------------------------------------------- %

\usetheme{geoslides}            % options: see the top of beamerthemegeoslides.sty
\usepackage{notation}           % maths shortcuts shared with the paper (optional)
\usepackage[main=british]{babel}

% Any other package the deck needs goes here, e.g. \usetikzlibrary{arrows.meta}.

% Paths are relative to this file: in a project with a shared bibliography,
% write e.g. \addbibresource{../references/references.bib}.
\addbibresource{references/references.bib}
\graphicspath{{figures/}}

% A logo at the right end of every frame title, and one on the title page:
% \logo{\includegraphics[height=0.8cm]{logos/nhh_logo_1f_positiv_helblaa}}
% \titlegraphic{\includegraphics[height=1.2cm]{logos/nhh_logo_1f_positiv_helblaa}}

% ---------------------------------------------------------------------------- %
%   Title page                                                                 %
% ---------------------------------------------------------------------------- %

\title[Short Title]{A Very Interesting Title}   % [short title for the footer]
\subtitle{A very interesting (optional) subtitle}

% One affiliation per author: each name is set with its affiliation under it.
\author{First Author \and Second Author}
\institute{University of the Black Metal \and University of the Shoegaze}

\date[Conference Name]{%
  Department of Whatever \\[0.1em]
  Institution \\[1em]
  1 January 2023%
}

% ---------------------------------------------------------------------------- %
%   Presentation                                                               %
% ---------------------------------------------------------------------------- %

\begin{document}

\maketitle

\begin{frame}{Outline}
  \tableofcontents
\end{frame}

% ---------------------------------------------------------------------------- %
\section{Introduction}
% ---------------------------------------------------------------------------- %

\begin{frame}{Outline}
  \tableofcontents[currentsection]
\end{frame}

% A frame label is what \slidelink and \slideback point at.
\begin{frame}[label=introduction]{Introduction}{A smart subtitle}

  \citet[p.~385]{bib:becker1993}:

  \begin{displayquote}
    ``The economic approach is a method of analysis, not an assumption of
    particular motivations. [\ldots] The analysis assumes that individuals
    maximize welfare as they conceive it, whether they be selfish,
    altruistic, loyal, spiteful, or masochistic.''
  \end{displayquote}

  \begin{itemize}
    \item A bullet with a link to a backup slide:
      \slidelink{derivation}{the derivation}.
    \item<2-> Overlays work as usual: this bullet appears on click.
    \item[$\rightarrow$]<2-> Custom markers too.
  \end{itemize}

\end{frame}

% One sentence, full bleed, on the inverted palette.
\begin{frame}[standout]
  The one sentence the room should leave with.
\end{frame}

% ---------------------------------------------------------------------------- %
\section{Results}
% ---------------------------------------------------------------------------- %

\begin{frame}{A Figure}{What it shows}
  \begin{figure}
    \includegraphics[width=.7\linewidth]{example-image}
    \caption{A caption in the house style: the long sentence that says
    what the reader is looking at.}
  \end{figure}
\end{frame}

% A table fragment written by the estimation code, as it comes.
\begin{frame}{A Table}{Estimates}
  \begin{table}
    \begin{threeparttable}
      % A stand-in for the table fragments an estimation script writes: main.tex
% pulls it in with \input.
\begin{tabular}{lccc}
  \toprule
  {}                 & (1)     & (2)     & (3)     \\
  {}                 & OLS     & OLS     & IV      \\
  \midrule
  Treatment          & 0.184   & 0.161   & 0.229   \\
                     & (0.042) & (0.045) & (0.081) \\
  Controls           & No      & Yes     & Yes     \\
  \midrule
  Observations       & 12,480  & 12,480  & 12,480  \\
  \bottomrule
\end{tabular}

      \begin{tablenotes}\scriptsize
      \item Notes: standard errors in parentheses.
      \end{tablenotes}
    \end{threeparttable}
    \caption{Difference-in-differences estimates of the outcome on
    treatment, with and without controls.}
  \end{table}
\end{frame}

\begin{frame}{Two Columns}
  \begin{columns}[T]
    \begin{column}{.46\textwidth}
      \begin{itemize}
        \item The claim on the left.
        \item The qualification under it.
      \end{itemize}
    \end{column}
    \begin{column}{.5\textwidth}
      \begin{figure}
        \includegraphics[width=\linewidth]{example-image}
        \caption{The evidence on the right.}
      \end{figure}
    \end{column}
  \end{columns}
\end{frame}

\begin{frame}{The Model}
  \begin{assumption}[Preferences]
    Utility is increasing and concave in consumption.
  \end{assumption}

  \begin{theorem}
    Under the assumption, the agent solves
    \begin{equation*}
      \max_{c \gs 0} \; \bE\squares{\sum_{t=0}^{\infty} \beta^t u(c_t)}.
    \end{equation*}
  \end{theorem}

  \begin{block}{Intuition}
    Blocks, \alert{alerts} and theorems take the palette's colours.
  \end{block}
\end{frame}

% ---------------------------------------------------------------------------- %
%   Closing and references                                                     %
% ---------------------------------------------------------------------------- %

\begin{frame}
  \centering
  {\LARGE\inserttitle\par}
  \vspace{3em}
  {\Large Thank you!\par}
  \vspace{1em}
  {\Large\href{mailto:your@email.com}{your@email.com}\par}
\end{frame}

% References are listed by author; citations stay in date order.
\begin{frame}[allowframebreaks]{References}
  \newrefcontext[sorting=nyt]
  \printbibliography[heading=none]
\end{frame}

% ---------------------------------------------------------------------------- %
%   Backup slides                                                              %
% ---------------------------------------------------------------------------- %

% \appendix keeps the backup slides out of the frame count. Its first frame is
% where the "Appendix" link in the footer leads: here, an index of the backup
% slides by section.
\appendix

\begin{frame}{Backup Slides}
  \tableofcontents
\end{frame}

\section{Derivations}

\begin{frame}[label=derivation]{A Backup Slide}{\slideback{introduction}}
  \begin{itemize}
    \item The detail the question was about.
  \end{itemize}
\end{frame}

\end{document}
